\documentclass[10pt]{article}
\usepackage[margin=0.85in]{geometry}
\usepackage{amsmath,amssymb,amsthm,mathtools}
\usepackage{microtype}
\usepackage{enumitem}
\usepackage{booktabs}
\usepackage{multicol}
\usepackage[hidelinks]{hyperref}
\setlist{nosep}
\newtheorem{theorem}{Theorem}[section]
\newtheorem{proposition}[theorem]{Proposition}
\newtheorem{corollary}[theorem]{Corollary}
\newtheorem{lemma}[theorem]{Lemma}
\theoremstyle{definition}
\newtheorem{definition}[theorem]{Definition}
\newtheorem{warning}[theorem]{Warning}
\theoremstyle{remark}
\newtheorem{remark}[theorem]{Remark}
\newcommand{\csch}{\operatorname{csch}}
\newcommand{\op}{\mathrm{op}}
\title{Proposed Manuscript Insertions IV\\
Full Trace-Response Plumbing and Exact Equipotential Collapse}
\author{Research draft for the marked-width MLC program}
\date{July 25, 2026}

\begin{document}
\maketitle

\begin{abstract}
A long planar plumbing neck is not, in general, a single scalar resistor:
nonconstant boundary traces carry a nonvanishing half-cylinder response. We
correct the finite-edge approximation by computing the complete
Dirichlet-to-Neumann trace operator on the canonical ray/equipotential
rectangle. Serial collars form an exact response semigroup, the full operator
is determined by total resistance and a universal tangential operator, and its
trace-renormalized nodal approximation has sharp relative error
$2/(e^{\pi R}-1)$. The estimate survives arbitrary Hilbert-space Schur
complements and preserves strict response contraction after an explicit finite
depth choice. Exact scalar collapse is characterized: it holds precisely when
all admissible traces are constant. These results close the direct-decoration
end-plumbing trace-response gate but do not prove the global response theorem,
MBTT, rigidity, or MLC.
\end{abstract}

\section{Correction to the scalar plumbing model}

Let
\[
 C_R=(0,R)\times(0,1)
\]
with natural Neumann conditions on the side walls. Prescribe
\[
 f(y)=\sqrt2\cos(\pi y),\qquad g(y)=0
\]
on the two end cross-sections. Both traces have mean zero, while their harmonic
extension has energy
\[
 \mathcal E_R(f,0)=\pi\coth(\pi R)\longrightarrow\pi.
\tag{1.1}
\]
A model retaining only the two trace averages assigns zero energy. Thus a
finite conductance edge represents the continuum neck only after the
cross-sections have been electrically shorted, or after all admissible traces
have separately been proved constant.

\begin{warning}[Correction of the previous nodal estimate]
The node \path{ADAPTIVE_NODAL_RESPONSE_APPROXIMATION} from the preceding
research insertion must be replaced. Its single-edge estimate is valid under
an explicit equipotential-electrode hypothesis. For arbitrary surface traces,
the full operator below is required.
\end{warning}

\section{Canonical trace collars}

Let $\Phi$ be the Bottcher coordinate in the basin of infinity, and write
\[
 G(z)=\log|\Phi(z)|,\qquad
 \Theta(z)=\frac{1}{2\pi}\arg\Phi(z)\pmod1.
\]
A sector bounded by rays $\Theta=\theta_0,\theta_1$ and equipotentials
$G=t_-,t_+$ is mapped to a width-one rectangle by
\[
 (x,y)=\left(
 \frac{G-t_-}{2\pi(\theta_1-\theta_0)},
 \frac{\Theta-\theta_0}{\theta_1-\theta_0}
 \right).
\]
Its length is
\[
 R=\frac{t_+-t_-}{2\pi(\theta_1-\theta_0)}.
\tag{2.1}
\]
For a univalent branch of $f^q$,
\[
 G(f^qz)=2^qG(z),\qquad
 \Theta(f^qz)=2^q\Theta(z)\pmod1.
\tag{2.2}
\]
Hence inverse branches scale potential thickness and angular width by the same
factor, preserving the normalized trace coordinate.

\begin{theorem}[Trace-compatible Bottcher collar]
Choose a parent equipotential cross-section $\gamma$ and put
$\gamma_w=\phi_w(\gamma)$. Let $C_{w,a}$ be the full external sector between
$\gamma_w$ and $\gamma_{wa}$, with the two authorized ray sides retained as
walls. Then $C_{w,a}$ is marked-conformally a rectangle of width one and some
length $R_a>0$. Adjacent child and parent trace parameterizations agree. Thus a
tail $a_1\cdots a_m$ is marked-conformally a rectangle of length
\[
 R(a_1\cdots a_m)=\sum_{j=1}^mR_{a_j}.
\tag{2.3}
\]
The construction is natural under every certified univalent presentation and
root enlargement that retains the collar.
\end{theorem}

\begin{remark}[Source-facing check]
A claimed application must verify that no protected collar, obstacle, or
localization boundary cuts through the selected external sector. Failure is
recorded at the first affected cell.
\end{remark}

\section{Exact full trace response}

Let
\[
 \varphi_0(y)=1,\qquad
 \varphi_n(y)=\sqrt2\cos(n\pi y),\quad n\ge1,
\]
be the Neumann eigenbasis on $(0,1)$. Put
\[
 A=\sqrt{-\partial_y^2{}_N},\qquad
 A\varphi_0=0,\quad A\varphi_n=n\pi\varphi_n,
\tag{3.1}
\]
and let $P_0$ be projection onto constants. The normalized trace space is
$\mathcal H=H^{1/2}_N(0,1)$ with
\[
 \|f\|_{\mathcal H}^2=|f_0|^2+\sum_{n\ge1}n\pi|f_n|^2.
\tag{3.2}
\]

\begin{theorem}[Full rectangle Dirichlet-to-Neumann formula]
For end traces $f,g\in\mathcal H$, the minimizing energy on $C_R$ is
\[
\begin{aligned}
 \mathcal E_R(f,g)
 ={}&\frac{|f_0-g_0|^2}{R}\\
 &+\sum_{n\ge1}n\pi\left[
 \coth(n\pi R)(|f_n|^2+|g_n|^2)
 -2\csch(n\pi R)\Re(f_n\overline{g_n})
 \right].
\end{aligned}
\tag{3.3}
\]
Equivalently,
\[
 \mathcal N_R=
 \begin{pmatrix}
 \alpha_R(A)&-\beta_R(A)\\
 -\beta_R(A)&\alpha_R(A)
 \end{pmatrix},
\tag{3.4}
\]
where
\[
 \alpha_R(\lambda)=\lambda\coth(R\lambda),\qquad
 \beta_R(\lambda)=\lambda\csch(R\lambda),
\tag{3.5}
\]
and both functions take the value $R^{-1}$ at $\lambda=0$.
\end{theorem}

\begin{proof}
For $n=0$, the minimizing mode is linear in $x$. For $n\ge1$, it is
\[
 u_n(x)=
 \frac{\sinh(n\pi(R-x))}{\sinh(n\pi R)}f_n
 +\frac{\sinh(n\pi x)}{\sinh(n\pi R)}g_n.
\]
Green's identity gives the displayed mode energy. Orthogonality completes the
sum.
\end{proof}

\begin{proposition}[Marked phase]
If the right trace is identified with the left normalized coordinate by a
unitary $U$ commuting with $A$ and preserving constants, then
\[
 \mathcal N_{R,U}=
 \begin{pmatrix}
 \alpha_R(A)&-\beta_R(A)U\\
 -U^*\beta_R(A)&U^*\alpha_R(A)U
 \end{pmatrix}.
\tag{3.6}
\]
Serial gluing multiplies the phases. For $R$ bounded below, the family over a
compact marked phase group is norm compact; the phase-dependent off-diagonal
part is smoothing.
\end{proposition}

\begin{theorem}[Exact response semigroup]
If two trace-compatible collars of lengths $R,S$ are glued along their common
parameterized cross-section, then
\[
 \boxed{\mathcal N_R\star\mathcal N_S=\mathcal N_{R+S},}
\tag{3.7}
\]
where $\star$ means addition of energy forms followed by minimization over the
common trace.
\end{theorem}

\begin{proof}
The glued surface is the rectangle of length $R+S$. Restricting its harmonic
minimizer gives the minimizing interface trace, and conversely the two
minimizers glue to an admissible function on the long rectangle.
\end{proof}

\begin{corollary}[Finite encoding of the full operator]
The complete trace response of a serial tail is determined exactly by its total
resistance $R$, the fixed tangential operator $A$, and the marked
phase/orientation. It is not an arbitrary infinite-dimensional state.
\end{corollary}

\section{Exact equipotential collapse}

Restricting (3.3) to constant traces gives
\[
 \mathcal E_R(a,b)=\frac{|a-b|^2}{R}.
\tag{4.1}
\]
Thus the constant mode is exactly a scalar resistor of conductance $1/R$.

\begin{theorem}[Exact mean-collapse criterion]
Let $V\subset\mathcal H$ be a linear admissible trace subspace and let
$M(f)=f_0$ be the trace-mean map. The restriction of $\mathcal E_R$ to
$V\oplus V$ factors through the two means if and only if $M|_V$ is injective.
If constants are admissible, $\mathbb C\varphi_0\subset V$, or more generally
if $f-M(f)\varphi_0\in V$ for every $f\in V$, this is equivalent to
\[
 V=\mathbb C\varphi_0=\ker A.
\tag{4.2}
\]
\end{theorem}

\begin{proof}
If the response factors through means and $M(f)=0$, then
$\mathcal E_R(f,0)=\mathcal E_R(0,0)=0$. Formula (3.3) forces $f=0$, so
$M|_V$ is injective. Conversely, injectivity gives a unique representative
$f_a\in V$ of each $a\in M(V)$; define
$q_R(a,b)=\mathcal E_R(f_a,f_b)$. If constants are admissible, then
$f-M(f)\varphi_0\in V\cap\ker M$, hence it vanishes. The same argument uses
only mean-subtraction invariance.
\end{proof}

\begin{remark}[One-dimensional shape exception]
Without constant admissibility or mean-subtraction invariance, a nonconstant
one-dimensional trace space can have injective mean map and therefore admit a
one-scalar encoding. This records one fixed nonconstant shape; it is not an
electrically shorted equipotential.
\end{remark}

\begin{corollary}[Shorted cross-section]
If a cross-section is declared one perfect electrical electrode, every
admissible trace is constant and the one-edge model is exact.
\end{corollary}

\begin{warning}
A geometric Bottcher equipotential is not automatically an electrical
equipotential for the response problem.
\end{warning}

\section{Trace-renormalized nodal response}

The mean-zero response of a semi-infinite rectangle is $A$. Therefore define
\[
 \mathcal E_R^{\#}(f,g)
 =\frac{|f_0-g_0|^2}{R}
 +\langle Af,f\rangle+\langle Ag,g\rangle,
\tag{5.1}
\]
with operator
\[
 \mathcal N_R^{\#}=
 \begin{pmatrix}
 A+R^{-1}P_0&-R^{-1}P_0\\
 -R^{-1}P_0&A+R^{-1}P_0
 \end{pmatrix}.
\tag{5.2}
\]
This is the correct nodal normalization: two semi-infinite trace ends plus the
exact resistor joining their constant modes.

\begin{theorem}[Sharp full trace estimate]
Put
\[
 \eta(R)=\frac{2}{e^{\pi R}-1}.
\tag{5.3}
\]
Then
\[
 \left|\mathcal E_R-\mathcal E_R^{\#}\right|
 \le\eta(R)
 \left(\langle Af,f\rangle+\langle Ag,g\rangle\right).
\tag{5.4}
\]
Whenever $\eta(R)<1$,
\[
 \boxed{
 (1-\eta(R))\mathcal N_R^{\#}
 \preceq\mathcal N_R
 \preceq(1+\eta(R))\mathcal N_R^{\#}.}
\tag{5.5}
\]
The modewise constant is sharp.
\end{theorem}

\begin{proof}
For $t=n\pi R$, the normalized difference matrix is
\[
 \begin{pmatrix}
 \coth t-1&-\csch t\\
 -\csch t&\coth t-1
 \end{pmatrix}.
\]
Its eigenvalues are
\[
 \frac{2}{e^t-1},\qquad -\frac{2}{e^t+1}.
\]
The largest absolute value is maximized at $n=1$. The constant modes agree
exactly.
\end{proof}

\begin{corollary}[Explicit depth]
To obtain relative trace error at most $\delta\in(0,1)$, it suffices that
\[
 R\ge\frac1\pi\log\left(1+\frac2\delta\right).
\tag{5.6}
\]
If every chosen cell has resistance at least $R_*>0$, it suffices to take
\[
 \ell\ge
 \left\lceil
 \frac1{\pi R_*}\log\left(1+\frac2\delta\right)
 \right\rceil.
\tag{5.7}
\]
\end{corollary}

\section{Compact operator limit}

Let
\[
 \mathcal N_\infty=
 \begin{pmatrix}A&0\\0&A\end{pmatrix},
\tag{6.1}
\]
with $A$ vanishing on constants.

\begin{theorem}[Full trace compactification]
In the normalized $\mathcal H\oplus\mathcal H$ to dual operator norm,
\[
 \|\mathcal N_R-\mathcal N_\infty\|_{\op}
 \le
 \max\left\{\frac2R,\frac{2}{e^{\pi R}-1}\right\}.
\tag{6.2}
\]
Hence $R\mapsto\mathcal N_R$ extends continuously to $R=\infty$.
\end{theorem}

\begin{proof}
The constant block has eigenvalues $0,2/R$. The mean-zero block is controlled
by the preceding theorem.
\end{proof}

\begin{remark}
The end tail therefore contributes a compact one-parameter family of complete
trace operators. The root-side joint trace response is transported exactly;
no new arbitrary operator complexity is created by the address depth.
\end{remark}

\section{Hilbert-space gluing and external response}

Let $Q(x,\phi)\ge0$ be the component energy for external terminal values
$x$ and interface trace $\phi$. For a neck form $N$, define
\[
 q_N(x)=\inf_\phi\{Q(x,\phi)+N(\phi)\}.
\tag{7.1}
\]

\begin{lemma}[Relative form bounds survive elimination]
If
\[
 (1-\delta)N_0\preceq N_1\preceq(1+\delta)N_0,
 \qquad0\le\delta<1,
\]
then
\[
 (1-\delta)q_{N_0}\le q_{N_1}\le(1+\delta)q_{N_0}.
\tag{7.2}
\]
\end{lemma}

\begin{proof}
Add the same nonnegative form $Q$ and take the infimum over the same trace
space.
\end{proof}

For disjoint tails $R_1,\ldots,R_s$, put
\[
 \eta_*=\max_i\frac{2}{e^{\pi R_i}-1}.
\tag{7.3}
\]

\begin{theorem}[Physical versus trace-renormalized response]
Let $\Lambda_{\mathrm{phys}}$ be the external response obtained with the exact
full necks, and let $\Lambda_{\#}$ be obtained with the trace-renormalized neck
forms. On the quotient by global constants,
\[
 \boxed{
 (1-\eta_*)\Lambda_{\#}
 \preceq\Lambda_{\mathrm{phys}}
 \preceq(1+\eta_*)\Lambda_{\#}.}
\tag{7.4}
\]
The same estimate holds after every fixed finite-electrode compression.
\end{theorem}

\begin{proof}
Sum the neck estimates, add the component energies, and apply the lemma.
Compression by an electrode map preserves Loewner order.
\end{proof}

\section{Preservation of a strict transition}

Let $L$ be a nonnegative, positively homogeneous, Loewner-monotone gauge. Then
(7.4) implies
\[
 (1-\eta_*)L(\Lambda_{\#})
 \le L(\Lambda_{\mathrm{phys}})
 \le(1+\eta_*)L(\Lambda_{\#}).
\tag{8.1}
\]

\begin{theorem}[Strict response transition survives]
Suppose the trace-renormalized normalized block satisfies
\[
 L_{\#}^{\mathrm{out}}
 \le\mu L_{\#}^{\mathrm{in}}+E,
 \qquad0<\mu<1.
\tag{8.2}
\]
If both endpoint materializations have trace error at most $\eta$, then
\[
 L_{\mathrm{phys}}^{\mathrm{out}}
 \le
 \mu\frac{1+\eta}{1-\eta}
 L_{\mathrm{phys}}^{\mathrm{in}}
 +(1+\eta)E.
\tag{8.3}
\]
The homogeneous factor remains strictly below one whenever
\[
 \eta<\frac{1-\mu}{1+\mu}.
\tag{8.4}
\]
Equivalently, a backward factor $\lambda=\mu^{-1}>1$ becomes
\[
 \lambda_{\mathrm{eff}}
 =\lambda\frac{1-\eta}{1+\eta}>1.
\tag{8.5}
\]
\end{theorem}

\begin{remark}[Exact retention]
If the state retains $\mathcal N_R$ through the exact functional calculus
(3.4), the plumbing approximation error is zero. The relative estimate is
needed only when passing to the compact trace-renormalized nodal model.
\end{remark}

\section{Promotion statement}

\begin{theorem}[End plumbing trace-response control]
For every source-authorized direct-decoration tail whose ray/equipotential
collar avoids all protected data and remains in the localization:
\begin{enumerate}
 \item the complete full trace response is exactly $\mathcal N_R$;
 \item serial gluing is exact and depends only on total resistance and marked
 phase;
 \item the operator family is compact at infinite depth;
 \item the trace-renormalized response has the sharp relative error (5.3);
 \item arbitrary component elimination and finite-electrode compression retain
 the same comparison;
 \item strict response contraction survives after the explicit depth choice;
 and
 \item scalar collapse is invoked only under the exact constant-trace
 criterion.
\end{enumerate}
Thus \texttt{END\_PLUMBING\_TRACE\_RESPONSE\_CONTROL} is proved in the
direct-decoration trace-compatible category.
\end{theorem}

\begin{remark}[Remaining global gate]
The theorem does not supply a trace collar for every mixed virtual-Molecule
block, a universally bounded response/WAD/Value quotient, a common KKL-neutral-
virtual chart, MBTT, coercivity, rigidity, or MLC. The broad node
\texttt{H2\_RESPONSE\_REALIZATION} remains open.
\end{remark}

\section{Ledger correction and next theorem}
\enlargethispage{8\baselineskip}

The previous single-edge node should be retained only as
\texttt{SHORTED\_EQUIPOTENTIAL\_NODAL\_RESPONSE}. New proved nodes are:
\begin{multicols}{2}
\small
\begin{enumerate}
 \item \texttt{BOTTCHER\_TRACE\_COLLAR\_ATLAS};
 \item \texttt{FULL\_TRACE\_RECTANGLE\_DTN};
 \item \texttt{TRACE\_RESPONSE\_SEMIGROUP};
 \item \texttt{EQUIPOTENTIAL\_COLLAPSE\_CRITERION};
 \item \texttt{TRACE\_RENORMALIZED\_NECK\_ESTIMATE};
 \item \texttt{END\_TRACE\_RESPONSE\_COMPACTIFICATION};
 \item \texttt{HILBERT\_TRACE\_GLUE\_CALCULUS};
 \item \texttt{TRACE\_STABLE\_STRICT\_TRANSITION\_TRANSFER}; and
 \item \texttt{END\_PLUMBING\_TRACE\_RESPONSE\_CONTROL}.
\end{enumerate}
\end{multicols}

The next source-facing target is a universal trace-interface authorization and
bounded future response theorem: every mixed block must supply its protected
trace collars, exact component joint response forms, enlargement maps, a
future-complete terminal/Value quotient, and one calibrated chart in which the
remaining non-plumbing error stays below the strict transition margin.

\end{document}
